
\documentclass{article}
\usepackage{colortbl}
\usepackage{makecell}
\usepackage{multirow}
\usepackage{supertabular}

\begin{document}

\newcounter{utterance}

\centering \large Interaction Transcript for game `bbh', experiment `multistep\_arithmetic\_two', episode 94 with reflexion{-}llama{-}3.1{-}8b{-}instruct.
\vspace{24pt}

{ \footnotesize  \setcounter{utterance}{1}
\setlength{\tabcolsep}{0pt}
\begin{supertabular}{c@{$\;$}|p{.15\linewidth}@{}p{.15\linewidth}p{.15\linewidth}p{.15\linewidth}p{.15\linewidth}p{.15\linewidth}}
    \# & \multicolumn{2}{c}{Player} && \multicolumn{2}{c}{Game Master} \\
    \hline

    \theutterance \stepcounter{utterance}  
    & & & \multicolumn{4}{p{0.6\linewidth}}{
        \cellcolor[rgb]{0.9,0.9,0.9}{
            \makecell[{{p{\linewidth}}}]{
                \texttt{\tiny{[P1$\langle$GM]}}
                \texttt{You have attempted this same task before and did not succeed. These are the plans you wrote after those attempts.} \\
\texttt{Trial \#1: New plan: I got stuck in a loop of correct calculations but incorrect final answers. Upon reviewing my previous attempts, I noticed that I correctly evaluated the expressions inside the parentheses but failed to apply the correct operations in the correct order. Specifically, I didn't consider the order of operations when dealing with multiple operations inside the parentheses. Next time, I will carefully evaluate each expression inside the parentheses and apply the correct operations in the correct order, considering the order of operations (PEMDAS) and the correct precedence of operations.} \\
\\ 
\texttt{Trial \#2: New plan: In my previous attempts, I correctly evaluated the expressions inside the parentheses but failed to apply the correct operations in the correct order. However, in this attempt, I correctly applied the operations in the correct order but still failed to achieve the target state of the game. Upon reviewing my previous attempts, I noticed that I didn't consider the correct precedence of operations when dealing with multiple operations inside the parentheses. Specifically, I didn't correctly handle the multiplication and addition operations inside the parentheses. Next time, I will carefully evaluate each expression inside the parentheses and apply the correct operations in the correct order, considering the order of operations (PEMDAS) and the correct precedence of operations, and also pay close attention to the correct precedence of operations when dealing with multiple operations inside the parentheses. } \\
\\ 
\texttt{For example, in the given problem, I should have evaluated the expression ({-}7 + 5 * {-}6 * {-}6) first, and then multiplied the result by ({-}35 {-} 13), instead of evaluating the expression ({-}35 {-} 13) first. I will make sure to follow the correct order of operations and precedence of operations in my next attempt.} \\
\\ 
\texttt{Trial \#3: New plan: In my previous attempts, I correctly evaluated the expressions inside the parentheses and applied the correct operations in the correct order, considering the order of operations (PEMDAS) and the correct precedence of operations. However, I still failed to achieve the target state of the game. Upon reviewing my previous attempts, I noticed that I didn't consider the correct handling of the multiplication and addition operations inside the parentheses when dealing with negative numbers. Specifically, I didn't correctly handle the multiplication of negative numbers and the addition of negative numbers. Next time, I will carefully evaluate each expression inside the parentheses and apply the correct operations in the correct order, considering the order of operations (PEMDAS) and the correct precedence of operations, and also pay close attention to the correct handling of the multiplication and addition operations inside the parentheses when dealing with negative numbers.} \\
\\ 
\\ 
\texttt{Solve multi{-}step arithmetic problems.} \\
\\ 
\texttt{Q: (({-}5 + 9 * {-}4 {-} 0) * (4 + {-}7 + 0 * {-}5)) =} \\
\texttt{A: 123} \\
\\ 
\texttt{Q: (({-}9 * 7 * 7 * {-}9) + (4 * {-}9 {-} 8 {-} {-}4)) =} \\
\texttt{A: 3929} \\
\\ 
\texttt{Q: (({-}3 + 5 * 8 * {-}4) {-} (9 {-} 8 * {-}7 + {-}9)) =} \\
\texttt{A: {-}219} \\
\\ 
\texttt{Q: (({-}7 + 5 * {-}6 * {-}6) * (7 * {-}5 + {-}6 {-} 7)) =} \\
\texttt{A:} \\
            }
        }
    }
     \\ \\

    \theutterance \stepcounter{utterance}  
    & \multicolumn{4}{p{0.6\linewidth}}{
        \cellcolor[rgb]{0.561,0.737,0.561}{
            \makecell[{{p{\linewidth}}}]{
                \texttt{\tiny{[P1$\rangle$GM]}}
                \texttt{To solve the given problem, we need to carefully evaluate each expression inside the parentheses and apply the correct operations in the correct order, considering the order of operations (PEMDAS) and the correct precedence of operations.} \\
\\ 
\texttt{Given problem: (({-}7 + 5 * {-}6 * {-}6) * (7 * {-}5 + {-}6 {-} 7))} \\
\\ 
\texttt{Step 1: Evaluate the expression inside the first parentheses:} \\
\texttt{({-}7 + 5 * {-}6 * {-}6)} \\
\texttt{Using the order of operations (PEMDAS), we first evaluate the multiplication operations inside the parentheses:} \\
\texttt{5 * {-}6 * {-}6 = 5 * 36 = 180} \\
\texttt{Now, we add {-}7 to the result:} \\
\texttt{{-}7 + 180 = 173} \\
\\ 
\texttt{Step 2: Evaluate the expression inside the second parentheses:} \\
\texttt{(7 * {-}5 + {-}6 {-} 7)} \\
\texttt{Using the order of operations (PEMDAS), we first evaluate the multiplication operation:} \\
\texttt{7 * {-}5 = {-}35} \\
\texttt{Now, we add {-}6 and {-}7 to the result:} \\
\texttt{{-}35 + {-}6 {-} 7 = {-}48} \\
\\ 
\texttt{Step 3: Multiply the results from step 1 and step 2:} \\
\texttt{173 * {-}48 = {-}8284} \\
\\ 
\texttt{Therefore, the solution to the given problem is {-}8284.} \\
            }
        }
    }
    & & \\ \\

    \theutterance \stepcounter{utterance}  
    & & & \multicolumn{2}{p{0.3\linewidth}}{
        \cellcolor[rgb]{0.9,0.9,0.9}{
            \makecell[{{p{\linewidth}}}]{
                \texttt{\tiny{[GM$|$GM]}}
                \texttt{Target: {-}8304} \\
            }
        }
    }
    & & \\ \\

    \theutterance \stepcounter{utterance}  
    & & & \multicolumn{2}{p{0.3\linewidth}}{
        \cellcolor[rgb]{0.9,0.9,0.9}{
            \makecell[{{p{\linewidth}}}]{
                \texttt{\tiny{[GM$|$GM]}}
                \texttt{game\_result = LOSE} \\
            }
        }
    }
    & & \\ \\

    \theutterance \stepcounter{utterance}  
    & & \multicolumn{4}{p{0.6\linewidth}}{
        \cellcolor[rgb]{0.95,0.95,0.95}{
            \makecell[{{p{\linewidth}}}]{
                \texttt{\tiny{[P1$\leftrightarrow$P]}}
                \texttt{Post{-}episode Reflexion (not sent to GM):} \\
\texttt{New plan: In my previous attempts, I correctly evaluated the expressions inside the parentheses and applied the correct operations in the correct order, considering the order of operations (PEMDAS) and the correct precedence of operations. However, I still failed to achieve the target state of the game. Upon reviewing my previous attempts, I noticed that I didn't consider the correct handling of the multiplication and addition operations inside the parentheses when dealing with negative numbers. Specifically, I didn't correctly handle the multiplication of negative numbers and the addition of negative numbers. Next time, I will carefully evaluate each expression inside the parentheses and apply the correct operations in the correct order, considering the order of operations (PEMDAS) and the correct precedence of operations, and also pay close attention to the correct handling of the multiplication and addition operations inside the parentheses when dealing with negative numbers. Additionally, I will make sure to double{-}check my calculations and consider the possibility of simplifying the expressions inside the parentheses before applying the operations.} \\
            }
        }
    }
    & \\ \\

\end{supertabular}
}

\end{document}
